\documentclass{article}
% Uncomment the following line to allow the usage of graphics (.png, .jpg)
\usepackage[pdftex]{graphicx}
%
\usepackage{hyperref}
%\usepackage[latin1]{inputenc}
\usepackage[utf8]{inputenc}
\usepackage[french]{babel}
\usepackage[T1]{fontenc}
%
\begin{document}
%
\begin{center}
\includegraphics[angle=0,width=0.5\textwidth,keepaspectratio='true']{graphics/xcsoarswiftsplash.pdf}
\vskip 0.5cm
\includegraphics[angle=0,width=0.66\textwidth,keepaspectratio='true']{graphics/title.pdf}
\subsection*{Prise en main du logiciel}
\end{center}
%
\begin{flushright}
Mise à jour du \today \space(par Adrien Ott)
\end{flushright}
%
\pagebreak
\noindent
\textbf{\textit{Avant-propos}}\\
\newline
Ce document est un guide de prise en main du logiciel XCSoar et ne constitue pas un mode d'emploi exhaustif du logiciel. L'objectif de ce document est de permettre à un(e) utilisateur(trice) ne connaissant pas le logiciel de l'installer, le configurer et l'utiliser le plus rapidement et facilement possible.\\Le mode d'emploi complet faisant actuellement référence est en anglais et disponible dans chaque répertoire de livraison (voir paragraphe "liens divers").\\Si vous avez des commentaires ou des suggestions d'amélioration de ce guide, n'hésitez pas à les faire sur le forum d'XCSoar, dans la rubrique "French".

%
% Just an example to inlude the images from the "images" directory
%\marginpar{\parbox{4.0cm}{\includegraphics[angle=0,width=1.0\linewidth,keepaspectratio='true']{images/smiley.png}}}

\section{Introduction}
\noindent
XCSoar est un logiciel d'aide à la navigation en planeur : "XC" est l'abbréviation de "cross country" (vol sur la campagne), "Soar" fait référence au vol sans moteur (vol à voile, vol libre).\\
%
\newline
%
XCSoar est disponible sur plusieurs platformes matérielles (PDA, PNA, smartphones, Altair de Triadis) et plusieurs systèmes d'exploitation, dont Windows Mobile 5 et 6, et Android (version 1.6 et plus récentes).\\
%
\newline
%
XCSoar est développé par des volontaires bénévoles avec un grand souci de qualité et de performance, et est très largement configurable du point de vue de l'interface utilisateur.\\
Ce guide de prise en main suppose que vous utilisiez XSoar version 6.3 ou plus récente.\\
Pour une version 6.x.y :
\begin{itemize}
\item les mises à jour dites de maintenance sont les évolutions "y" (correctifs si nécessaire),
\item les mises à jour majeures sont les évolutions "x".
\end{itemize}
%
Pour être informé des mises à jour, vous pouvez vous abonner au flux RSS du site XCSoar.org (voir paragraphe \ref{LiensUtiles} "Liens utiles").\\
%
\section{Installation}
\noindent
La première étape nécessaire à l'installation est le téléchagement de la dernière version d'XCSoar au lien suivant : \url{http://www.xcsoar.org/download/}\\
Sur cette page figurent les différents systèmes d'exploitation supportés, sélectionnez le votre et téléchagez le logiciel. Copiez le sur votre appareil, et lancez-le.\\
%
\newline
%
\itshape Sous Windows Mobile, préférez le fichier ".exe" plutôt que le ".cab", quitte à placer un raccourci vers le ".exe" dans le menu "Démarrer" de Windows.\upshape\\
\newline
XCSoar affiche alors le logo au centre, la version du logiciel en haut à gauche, un bouton "Quit" en haut à droite, et deux boutons de lancement \textbf{FLY} et \textbf{SIM} en bas de l'écran, respectivement pour le mode vol (pour voler et enregistrer la trace du vol) et le mode simulation (pour apprendre à l'utiliser et le configurer, de façon isolée ou avec votre appareil connecté à un simulateur, comme Condor par exemple). A ce stade vous pouvez quitter le logiciel (bouton "Quit").\\
%
Si vous avez lancé le mode \textbf{FLY} ou \textbf{SIM}, le menu principal peut être invoqué par un double clic sur le fond de carte.
L'utilisation est détaillée plus bas.
%
\section{Configuration}
%
\subsection{Les fichiers}
\noindent
XCSoar stocke tous les fichiers et répertoires nécessaires à son fonctionnement dans un répertoire nommé XCSoarData. Si vous diposez d'une carte mémoire, ce répertoire doit être à la racine (typiquement "/sdcard/XCSoarData"). Sinon le répertoire XCSoarData sera dans "Mes Documents" pour les systèmes Windows (c'est la configuration par défaut qu'il peut être préférable d'éviter sur les appareils ayant  peu de mémoire interne).\\
Les fichiers les plus utiles sont les suivants :
\begin{itemize}
\item points de virage ("waypoints"), ils peuvent être de diffèrents formats (*.dat de WinPilot/Cambridge, *.cup de SeeYou, *.wpz de Zander, *.wpt d'OziExplorer, GPSDump/FS, GEO ou UTM),
\item espaces aériens (au format OpenAir),
\item fond cartographique (format .xcm).
\end{itemize}
%
Pour trouver ces fichiers, voir le paragraphe \ref{LiensUtiles} "Liens utiles".
%
\subsection{L'interface}
\noindent
%
Une fois ces fichiers placés dans le répertoire XCSoarData, la configuration proprement dite peut commencer. L'interaction avec le logiciel se fait principalement via l'écran tactile (plus les boutons éventuels, comme sur un iPaq ou un smartphone Android, mais cela ne sera pas détaillé ici).\\
%
\newline
%
Lancez XCSoar en mode simulation et lancez le menu principal:
\begin{itemize}
\item double clic sur la carte,
ou
\item raccourci gestuel "bas-haut" (\textbf{B H}).
\end{itemize}
Un raccourci gestuel ("gesture" dans le manuel en anglais), est un dessin géométrique que l'utilisateur peut tracer à l'écran pour invoquer un menu en particulier. "\textbf{B H}" consiste à tracer sur l'écran, avec le doigt, un aller-retour commençant par descendre. D'autres exemples sont détaillés plus bas.\\
%
\newline
%

Le \textbf{menu principal} est constitué des boutons
\begin{itemize}
\item "Nav." pour gérer les fonctions du logiciel liées à la navigation,
\item "Affich." pour gérer les options d'affichage,
\item "Config." pour la configuration du logiciel,
\item "Info." pour consulter voire modifier l'état du système,
\item "Annuler", pour effacer ces boutons,
\item "Quitter" pour quitter le logiciel.
\end{itemize}
%
Une fois le menu principal invoqué, vous pouvez changez la langue utilisée par l'interface si celle-ci n'est pas le français, par exemple (si en anglais) : choisir "fr.mo" dans Config / Config / System Setup / Look Language, Inputs / Language.\\
%
\newline
%
Sinon passez directement à la configuration du logiciel dans le menu Config. / Config. / Options Système / Fichiers / Fichiers, et sélectionnez les fichiers correspondants aux champs suivants :
\begin{itemize}
\item Carte,
\item Waypoints,
\item Espaces aériens.
\end{itemize}
Les champs configurables sont affichés sur fond blanc et modifiables en cliquant dessus, puis en sélectionnant l'option voulue. Pour chaque option une aide est disponible en français, fournissant les explications sur les différentes options (bouton "Aide").\\
%
\newline
%
XCSoar permet l'affichage de \textbf{plusieurs pages} (de une à huit) dont le contenu de chaque InfoBox peut être paramétré. Voir dans "Config. / Config. / Options Système / Apparence" les menus "Pages des InfoBoxes" et "Modes InfoBoxes".\\
Une page nommée AUTO permet de basculer automatiquement entre les trois suivantes :
\begin{itemize}
\item Thermique,
\item Transition,
\item Arrivée.
\end{itemize}
Rien ne vous empêche de ne pas utiliser la page AUTO, ou de l'utiliser en parametrant les mêmes InfoBox pour Transition et Arrivée.\\
A l'utilisation, le passage d'une page à l'autre s'effectue au choix par le menu principal (appuyez deux fois sur Affich., puis le bouton correspondant au nom de la prochaine page) ou plus simplement par les raccourcis gestuels suivants :
\begin{itemize}
\item "gauche" (glisser le doigt de la droite vers la gauche sur la carte),
\item "droite" (l'inverse).
\end{itemize}
Ce qui permet respectivement de passer à la page précédente ou suivante (telles que listées dans le menu de configuration "Apparence / Pages de InfoBox").\\
%
L'affichage portrait ou paysage est également configurable, ainsi que la disposition des InfoBox sur l'écran (quoique identique sur toutes les pages).\\
%
\newline
%
Les \textbf{InfoBox} sont des fenêtres d'information numériques, alphanumériques ou graphiques. Leur configuration peut se faire au choix 
\begin{itemize}
\item dans le menu de configuration de la page correspondante (Config. / Config. / Options Système / Apparence / Modes InfoBox)
ou
\item en agissant directement sur l'InfoBox, à partir de la carte : en effectuant un appui long sur l'InfoBox correspondante.
\end{itemize}
%
\textit{Conseils :
\begin{itemize}
\item La capacité de configuration est telle qu'il est préférable que vos InfoBox reflétent vos habitudes de vol, ainsi demandez-vous de quelles informations vous avez besoin aujourd'hui, en évitant (dans un premier temps) d'afficher des données que vous n'utilisez pas,
\item Si vous utilisez la même information sur deux pages (par exemple la distance au prochain point de virage), alors positionnez la au même endroit sur l'écran.
\end{itemize} }
%
\noindent
\newline
Une fonction utile est l'affichage sur la carte du \textbf{plan nécessaire} pour rejoindre un aérodrome par exemple (en prenant en compte la hauteur de sécurité configurée). Cette fonction est activable dans "Config. / Config. / Options Système /  Afficher la carte / Eléments / Hauteur d'arrivée", y choisir "Finesse nécessaire".
%
\subsection{Raccourcis gestuels}
\noindent
%
Voici quelques raccourcis facilitant l'accès à certaines fonctions fréquemment utilisées :
\begin{itemize}
\item Menu circuit : \textbf{D B}
\item Menu waypoints : \textbf{B D}
\item Dégagements : \textbf{B G}
\item Page "Etats" : \textbf{G B D B G} (décrire un "S" pour \textit{Status})
\end{itemize}
%
\section{Utilisation}
\noindent
Les valeurs numériques de certaines InfoBox sont modifiables :
\begin{itemize}
\item  En mode SIM (e.g. altitude, route instantanée) afin de modifier la trajectoire du symbole,
\item En mode FLY (e.g. l'InfoBox du calage MacCready).
\end{itemize}
\noindent
\newline
La flêche verticale à gauche de l'écran indique l'écart au plan vertical pour rejoindre le prochain point de virage avec le calage sélectionné. Une seconde flêche peut être ajoutée à l'intérieur de la première pour indiquer l'écart au plan avec un calage nul. L'apparition d'une croix sur le symbole d'écart au plan signifie que la trajectoire directe vers le point de virage intersecte le terrain (prenant en compte la marge paramétrable de franchissement du terrain).\\
\section{Liens Utiles}
\noindent
\label{LiensUtiles}
\subsection{Liens XCSoar}
\noindent
Site du logiciel XCSoar :\\
\url{http://www.xcsoar.org}\\
\newline
Flux RSS du site XCSoar.org :\\
\url{http://xcsoar.org/atom.xml}\\
\newline
Rubrique de téléchargement de la dernière version :\\
\url{http://www.xcsoar.org/download/}\\
\newline
Génération de cartes :\\
\url{http://mapgen.xcsoar.org/}\\
\newline
Boîte à bugs (et demandes d'amélioration) :\\ \url{http://www.xcsoar.org/trac/}\\
\newline
Le forum, dont une rubrique en français pour poser vos questions :\\
\url{http://www.xcsoar.org/forum/}\\
\newline
Pour les développeurs, voir le "developer manual".
%
\subsection{Liens divers}
\noindent
\label{LiensDivers}
Page brute contenant toutes les versions officielles :\\ \url{http://download.xcsoar.org/releases/}\\
\newline
Différents sites pour points de virages et terrains :\\
\url{http://planeur.net/index.php?option=com_jooget&Itemid=177}\\
\url{http://qfu.free.fr/pdv.php}\\
\url{http://www.segelflug.de/vereine/welt2000/}\\
\url{http://www.cumulus-soaring.com/soaring_links/airports.htm}\\
\url{http://soaringweb.org/TP}\\
\url{http://www.lvzc.be/index.php/secretariaat/documenten/cat_view/62-keerpuntbestanden}\\
\newline
Espaces aériens en France :\\
\url{http://www.ffvvespaceaerien.org/?page_id=412}\\
Et ailleurs :\\
\url{http://soaringweb.org/Airspace}\\
%
\end{document}
